\section{Finite Gaussian moments and whole-space Sobolev compactness}
\label{S:compact-gaussian}

We prove the finite Gaussian and compact-cutoff estimates needed to construct
the wave operator path. All traces and Schatten norms in this section are
unnormalized. The spatial states remain on $\R^d$, $d \in \{1,2,3\}$. Here $p$
is an auxiliary Schatten and Gaussian-moment index. The auxiliary periodic space
below is used only to factor a compactly supported multiplication operator. No
infinite Gaussian series or time path is constructed in this section.

\subsection{Finite Gaussian matrices with both variance powers}

\LeanMarker{L:gaussian-scalar-powers}
\begin{lemma}[Scalar power estimates]\label{L:gaussian-scalar-powers}
  For $0 < \theta < 1$ and $u,v \ge 0$,
  \begin{equation}\label{E:gaussian-weighted-power}
    u^\theta v^{1-\theta} \le \theta u+(1-\theta)v.
  \end{equation}
  Consequently, if $L > 0$ and $0 \le q \le L$ are integers, then for real
  $x,y$,
  \begin{equation}\label{E:gaussian-integer-power}
    |x|^q|y|^{L-q}
    \le \frac qL|x|^L+\frac{L-q}L|y|^L.
  \end{equation}
  For nonnegative finite sequences, the usual H\"older inequality holds with
  exponents $m$ and $m/(m-1)$, $m > 1$. If $Z \ge 0$ has finite expectation $F$
  and $0 < r < 1$, then $\E Z^r \le F^r$.
\end{lemma}

\begin{proof}
  We first prove the weighted scalar inequality and its integer-power
  consequence. We then use normalization to obtain the two remaining bounds.

  \medskip
  \noindent\textit{Step 1.~}We prove the two scalar inequalities.
  If $v > 0$, divide~\eqref{E:gaussian-weighted-power} by $v$ and put $t = u/v$.
  The function $\theta t+1-\theta-t^\theta$ decreases to zero at one and then
  increases, by its derivative $\theta(1-t^{\theta-1})$. Continuity covers $t =
  0$; if $v = 0$ the inequality is immediate. For $0 < q < L$, apply this with
  $u = |x|^L$, $v = |y|^L$ and $\theta = q/L$. The endpoint cases $q = 0,L$ are
  identities, with the ordinary integer convention $a^0 = 1$.

  \medskip
  \noindent\textit{Step 2.~}We prove the finite sequence inequality.
  To prove finite H\"older, divide the two sequences by their respective
  $\ell^m$ and $\ell^{m/(m-1)}$ norms. If either norm vanishes, the result is
  immediate. Apply~\eqref{E:gaussian-weighted-power} with $\theta = 1/m$ to the
  corresponding normalized powers and sum; the sum of products is at most one.
  Undo the normalization.

  \medskip
  \noindent\textit{Step 3.~}We prove the expectation bound.
  Finally, $F = 0$ gives $Z = 0$ almost surely. If $F >
  0$, apply \eqref{E:gaussian-weighted-power} to $u = Z/F$, $v = 1$ and $\theta
  = r$, and integrate. The right side has expectation one. This proves
  Lemma~\ref{L:gaussian-scalar-powers}.
\end{proof}

\LeanMarker{L:gaussian-trace-bounds}
\begin{lemma}[Two Hermitian trace bounds]\label{L:gaussian-trace-bounds}
  Let $B,Y$ be finite Hermitian matrices, $m \ge 2$ an integer, and $L = 2m-2$.
  For integers $0 \le q \le L$,
  \begin{equation}\label{E:gaussian-trace-word}
    \operatorname{Re}\operatorname{Tr}(BY^qBY^{L-q})
    \le \operatorname{Tr}(B^2|Y|^L).
  \end{equation}
  If $D \ge 0$ is Hermitian, then
  \begin{equation}\label{E:gaussian-trace-holder}
    \operatorname{Tr}(D|Y|^{2m-2})
    \le (\operatorname{Tr}D^m)^{1/m}
    (\operatorname{Tr}|Y|^{2m})^{(m-1)/m}.
  \end{equation}
\end{lemma}

\begin{proof}
  We work in an eigenbasis of $Y$ throughout. The first bound follows by
  comparing scalar powers; the second follows by finite H\"older.

  \medskip
  \noindent\textit{Step 1.~}We bound the trace word.
  Diagonalize $Y$, with real eigenvalues $\lambda_i$, and write $b_{ij}$ for the
  entries of $B$ in that basis. Hermitian symmetry gives $b_{ji} =
  \overline{b_{ij}}$. The trace word is the real sum
  \[
    \sum_{i,j}|b_{ij}|^2\lambda_j^q\lambda_i^{L-q}.
  \]
  Bound each product by its absolute value and use
  \eqref{E:gaussian-integer-power}. After interchanging $i,j$, the two resulting
  sums agree because $|b_{ij}| = |b_{ji}|$. Each equals
  $\sum_{i,j}|b_{ij}|^2|\lambda_i|^L = \operatorname{Tr}(B^2|Y|^L)$. The two
  weights sum to one, proving \eqref{E:gaussian-trace-word}. No commutation of
  $B$ and $Y$ is used.

  \medskip
  \noindent\textit{Step 2.~}We bound the trace involving $D$.
  In the same basis put $d_i = D_{ii} \ge 0$. If $\theta_j$ are the nonnegative
  eigenvalues of $D$, a unitary change of basis gives $d_i =
  \sum_j|U_{ij}|^2\theta_j$. Both row and column sums of $|U_{ij}|^2$ equal one.
  Finite convexity therefore yields
  \[
    \sum_i d_i^m \le \sum_{i,j}|U_{ij}|^2\theta_j^m
    = \operatorname{Tr}D^m.
  \]
  The row sums provide the convex weights for each $d_i$. The column sums
  then give the equality after summing over $i$.
  For completeness, the convexity step follows from $t^m \ge s^m+ms^{m-1}(t-s)$
  for $s,t \ge 0$, obtained by differentiating the difference. Average this
  inequality with $s$ equal to the weighted mean. Now apply finite H\"older from
  Lemma~\ref{L:gaussian-scalar-powers} to $\sum_i d_i|\lambda_i|^{2m-2}$. Since
  $(2m-2)m/(m-1) = 2m$, the last display proves \eqref{E:gaussian-trace-holder},
  including zero sums. This proves Lemma~\ref{L:gaussian-trace-bounds}.
\end{proof}

The estimate below is a finite-dimensional Gaussian noncommutative Khintchine
inequality. For its origins in the work of Lust-Piquard and a Gaussian
integration-by-parts proof, see Tropp~\cite{tropp:source2018:second-order}*{arXiv:1504.05919v2,
Proposition~2.3 and Section~7}. We include the proof and Hermitian dilation
to fix the normalization and retain both variance operators.

\LeanMarker{P:finite-gaussian-matrix}
\begin{proposition}[Finite Gaussian matrix moment]\label{P:finite-gaussian-matrix}
  Let $A_1,\ldots,A_N$ be complex $o$-by-$n$ matrices, with $o,n,N \ge 1$, and
  let $g_k$ be independent real standard Gaussians. Put
  \[
    X=\sum_k g_kA_k,\qquad V_{\rm out}=\sum_k A_kA_k^*,
    \qquad V_{\rm in}=\sum_k A_k^*A_k.
  \]
  For even $p = 2m \ge 2$, with $\|A\|_{S_p}^p = \operatorname{Tr}((A^*A)^m)$,
  \begin{align}
    (\E\|X\|_{S_p}^p)^{1/p}
    &\le \sqrt{p-1}\left(
    \frac{\operatorname{Tr}V_{\rm out}^m+
    \operatorname{Tr}V_{\rm in}^m}{2}\right)^{1/p}
    \label{E:finite-gaussian-average}\\
    &\le \sqrt{p-1}\max\left(
    \|V_{\rm out}^{1/2}\|_{S_p},\|V_{\rm in}^{1/2}\|_{S_p}\right).
    \label{E:finite-gaussian-maximum}
  \end{align}
  The constants are independent of $o,n,N$. The first inequality is an equality
  at $p = 2$. Empty or zero families give the zero identity.
\end{proposition}

\begin{proof}
  We first prove the moment bound for Hermitian matrices by Gaussian
  integration by parts. We then use Hermitian dilation to treat rectangular
  matrices and keep track of the factor two.

  \medskip
  \noindent\textit{Step 1.~}We justify Gaussian integration by parts.
  First suppose $B_k$ are Hermitian and put $Y = \sum_k g_kB_k$ and $D = \sum_k
  B_k^2$. All polynomial moments are finite because Gaussian decay dominates
  every fixed power. Componentwise integration by parts against the product
  Gaussian density gives $\E[g_kP(g)] = \E[\partial_kP(g)]$, with zero boundary
  terms. It applies to complex polynomials through their real and imaginary
  parts; absolute polynomial integrability justifies Fubini.

  \medskip
  \noindent\textit{Step 2.~}We estimate the Hermitian moment.
  For $m \ge 2$, put $F = \E\operatorname{Tr}Y^{2m}$ and $L = 2m-2$. Write
  $\operatorname{Tr}Y^{2m} = \sum_k g_k \operatorname{Tr}(B_kY^{2m-1})$ and
  differentiate the $2m-1$ factors in their order. This gives exactly
  \[
    F=\sum_k\sum_{q=0}^{L}
    \E\operatorname{Tr}(B_kY^qB_kY^{L-q}).
  \]
  The eigenbasis calculation above shows these traces are real. The
  differentiated polynomial is real by conjugation and trace cyclicity. No
  derivative of a spectral absolute value or measurable choice of an eigenbasis
  is needed. The two trace bounds give
  \begin{align*}
    F&\le (2m-1)\E\operatorname{Tr}(D|Y|^{2m-2})\\
    &\le (2m-1)(\operatorname{Tr}D^m)^{1/m}
    \E(\operatorname{Tr}|Y|^{2m})^{(m-1)/m}\\
    &\le (2m-1)(\operatorname{Tr}D^m)^{1/m}F^{(m-1)/m}.
  \end{align*}
  The first inequality bounds each trace word and sums the $2m-1$
  positions of the derivative. The second applies the trace H\"older bound.
  Here $\operatorname{Tr}|Y|^{2m} = \operatorname{Tr}Y^{2m}$, and the last
  inequality is the expectation bound in Lemma~\ref{L:gaussian-scalar-powers}.
  If $F = 0$ the conclusion is immediate. Otherwise divide by its positive
  fractional power and raise to $m$ to obtain
  \begin{equation}\label{E:finite-hermitian-gaussian}
    \E\operatorname{Tr}Y^{2m} \le (2m-1)^m\operatorname{Tr}D^m.
  \end{equation}
  For $m = 1$, independence gives directly $\E\operatorname{Tr}Y^2 =
  \sum_k\operatorname{Tr}B_k^2 = \operatorname{Tr}D$, proving the same formula
  as equality.

  \medskip
  \noindent\textit{Step 3.~}We pass to rectangular matrices.
  For rectangular $A_k$, use the Hermitian dilation
  \[
    B_k=\begin{pmatrix}0&A_k\\A_k^*&0\end{pmatrix},\qquad
    Y=\begin{pmatrix}0&X\\X^*&0\end{pmatrix},\qquad
    D=\operatorname{diag}(V_{\rm out},V_{\rm in}).
  \]
  Finite entrywise summation proves $\operatorname{Tr}(UV) =
  \operatorname{Tr}(VU)$ for compatible rectangular matrices. Since $(XX^*)^m =
  X(X^*X)^{m-1}X^*$, it follows that
  \[
    \operatorname{Tr}Y^{2m}=2\operatorname{Tr}((X^*X)^m),\qquad
    \operatorname{Tr}D^m=\operatorname{Tr}V_{\rm out}^m+
    \operatorname{Tr}V_{\rm in}^m.
  \]
  Substitute in~\eqref{E:finite-hermitian-gaussian}, divide by two and take the
  $p$th root. This proves~\eqref{E:finite-gaussian-average}. For positive $V$,
  spectral diagonalization gives $\|V^{1/2}\|_{S_p}^p = \operatorname{Tr}V^m$,
  and the average of two nonnegative numbers is bounded by their maximum. At $p
  = 2$, $\E\operatorname{Tr}X^*X = \operatorname{Tr}V_{\rm in} =
  \operatorname{Tr}V_{\rm out}$, proving the claimed equality. This proves
  Proposition~\ref{P:finite-gaussian-matrix}.
\end{proof}

Both variance directions are needed: the row-matrix example in
No-go~\ref{P:nogo-23} rules out a dimension-independent estimate that
keeps only one of them.


\LeanMarker{R:gaussian-one-variance-test}
\begin{remark}[Both variance directions are necessary]
  \label{R:gaussian-one-variance-test}
  Let $n = 16$ and $A_k = e_k^{\mathsf{T}}$ be the coordinate row matrices. The
  only nonzero singular value of $X$ is $(\sum_k g_k^2)^{1/2}$. Gaussian
  integration by parts gives $\E g^2 = 1$, $\E g^4 = 3$, and hence
  \[
    \E\|X\|_{S_4}^4=3n+n(n-1)=n(n+2)=288.
  \]
  Here $V_{\rm out} = n$ and $V_{\rm in} = I_n$. Keeping only the input variance
  in~\eqref{E:finite-gaussian-maximum} would give fourth power
  $9\operatorname{Tr}I_n^2 = 144$, which is false. Transposing the family
  refutes the corresponding output-only bound. The average and maximum bounds
  have fourth powers $1224$ and $2304$, respectively. This finite example
  asserts nothing about a specialized wave estimate.
\end{remark}

\subsection{Whole-space cutoffs through an auxiliary periodic space}

Use the unitary Fourier transform with $c_d = (2\pi)^{-d/2}$ and $\Lambda =
(1-\Delta)^{1/2}$. For real $r$, the norm on $H^r(\R^d)$ is
$\|\langle\xi\rangle^r\widehat{u}\|_2$. The deterministic cutoff arguments allow
any smoothing order $\gamma > 0$. Fix real $\sigma$ and a real compact smooth
$\chi$. Singular values $s_j$, $j \ge 1$, are counted with multiplicity, with
zeros appended when necessary. For compact $K$, put $\|K\|_{S_p}^p = \sum_j
s_j(K)^p$.

\begin{lemma}[Cutoff factorization]\label{L:cutoff-factorization}
  Choose a cube $Q_L = (-L/2,L/2)^d$ containing $\operatorname{supp}\chi$ in its
  interior, and a real $\psi \in C_c^\infty(Q_L)$ equal to one near
  $\operatorname{supp}\chi$. On the auxiliary torus $\mathbb{T}_L$ use Lebesgue
  measure and the orthonormal basis
  \[
    e_k(x)=L^{-d/2}e^{i\kappa_k\cdot x},\qquad
    \kappa_k=2\pi k/L,\qquad k \in \mathbb{Z}^d,
  \]
  with $\|v\|_{H^r(\mathbb{T}_L)}^2 = \sum_k\langle\kappa_k\rangle^{2r}|v_k|^2$.
  Periodization of $\psi u$ and multiplication of a periodic function by $\chi$
  extend to bounded maps
  \begin{equation}\label{E:cutoff-periodic-maps}
    \mathcal{P}_\psi:H^r(\R^d)\to H^r(\mathbb{T}_L),\qquad
    \mathcal{E}_\chi:H^r(\mathbb{T}_L)\to H^r(\R^d)
  \end{equation}
  for every real $r$. If $\mathcal{I}_\gamma$ denotes the identity inclusion
  from $H^{\sigma+\gamma}(\mathbb{T}_L)$ to $H^\sigma(\mathbb{T}_L)$, then
  \begin{equation}\label{E:cutoff-whole-space-recovery}
    \mathcal{E}_\chi\mathcal{I}_\gamma\mathcal{P}_\psi u=\chi u,
    \qquad u \in H^{\sigma+\gamma}(\R^d).
  \end{equation}
\end{lemma}

\begin{proof}
  We have to show two things. First, the two maps
  in~\eqref{E:cutoff-periodic-maps} are bounded for every real $r$. Second, the
  identity~\eqref{E:cutoff-whole-space-recovery} holds on
  $H^{\sigma+\gamma}(\R^d)$. In both cases we argue first on a dense class,
  namely Schwartz functions on $\R^d$ and finite torus polynomials on
  $\mathbb{T}_L$, and then pass to the whole space by density. The proof has
  five steps. Step~1 collects three elementary facts. Step~2 recalls the mixed
  Schur bound. Step~3 writes the two maps as kernel operators in the Fourier
  variables and bounds them. Step~4 extends the maps by density. Step~5 proves
  the identity~\eqref{E:cutoff-whole-space-recovery}.

  \medskip
  \noindent\textit{Step 1.~}We collect three elementary facts. The first fact
  compares the weights at
  two points. For $x,y \in \R^d$ and real $r$, the elementary inequality
  $\langle x\rangle \le \sqrt{2}\langle y\rangle\langle x-y\rangle$ gives
  \begin{equation}\label{E:cutoff-peetre}
    \left(\frac{\langle x\rangle}{\langle y\rangle}\right)^r
    \le 2^{|r|/2}\langle x-y\rangle^{|r|}.
  \end{equation}
  For $r \ge 0$ we raise the elementary inequality to the power $r$. For
  negative $r$, we interchange $x,y$ before raising to the power $|r|$.

  The second fact is that integration by parts gives rapid decay of the Fourier
  transform of every compact smooth cutoff.

  The third fact is a lattice sum that is bounded uniformly in its shift. For
  any lattice spacing $a > 0$ and $N > d$,
  \begin{equation}\label{E:cutoff-lattice-sum}
    \sup_y\sum_{k\in\mathbb{Z}^d}\langle y-ak\rangle^{-N}<\infty.
  \end{equation}
  Indeed, the sum does not change when $y$ is shifted by a lattice vector. We
  translate $y$ into a centered lattice cell, where $|y| \le a\sqrt d/2$. For
  $|k| \ge \sqrt d$ the triangle inequality gives $|y-ak| \ge a|k|-|y|$, and
  we have $|y-ak| \ge a|k|/2$. We treat the terms with $|k| < \sqrt d$ and the
  tail $|k| \ge \sqrt d$ separately. The remaining terms, those with $|k| <
  \sqrt d$, are finite: each is at most one, and their number depends only on
  $d$. The tail converges by dyadic lattice counting, since there are at most
  $C_d2^{jd}$ lattice points in the $j$th shell $2^j \le |k| < 2^{j+1}$ in
  $k$. More precisely, for $|k| \ge \sqrt d$, the preceding inequality gives
  \[
    \langle y-ak\rangle^{-N}
    \le \max(1,2/a)^N\langle k\rangle^{-N}.
  \]
  Here we use $\langle y-ak\rangle \ge \langle ak/2\rangle \ge
  \min(1,a/2)\langle k\rangle$. The comparison-series terms $\langle
  k\rangle^{-N}$ are bounded by $2^{-jN}$ on each dyadic $k$-shell. Hence, the
  comparison series is bounded by the geometric sum $\sum_{j \ge 0}
  C_d2^{j(d-N)}$, which converges because $N > d$. Thus, the resulting bound
  carries constants depending on $a,d,N$, as
  does~\eqref{E:cutoff-lattice-sum}.

  \medskip
  \noindent\textit{Step 2.~}We recall the mixed Schur bound. We use the mixed
  Schur bound for a
  kernel whose absolute row and column integrals are bounded by $M_1,M_2$. It
  states that the squared operator norm of the kernel operator is at most
  $M_1M_2$. To prove it, we apply the weighted Cauchy--Schwarz inequality at
  each output point and then Tonelli's theorem. The same proof works with
  counting measure on one side and Lebesgue measure on the other. For counting
  measure inputs, use finitely supported sequences; for Lebesgue measure
  inputs, use smooth compactly supported functions. The estimate extends to
  the corresponding $L^2$ spaces by density. For the kernels below, the same
  Cauchy--Schwarz estimate on truncated inputs identifies the extension with
  the stated integral or sum.

  \medskip
  \noindent\textit{Step 3.~}We prove the bounds
  in~\eqref{E:cutoff-periodic-maps}. We write each map as a kernel operator
  between the weighted Fourier variables and apply Step~2. We begin with
  $\mathcal{P}_\psi$. For Schwartz $u$, the $k$th coefficient of its cutoff
  periodization is
  \[
    (\mathcal{P}_\psi u)_k
    = L^{-d/2}\int\psi(x)u(x)e^{-i\kappa_k\cdot x}\ud x
    = L^{-d/2}\int\widehat{\psi}(\kappa_k-\xi)\widehat{u}(\xi)\ud\xi.
  \]
  In the second equality, the $c_d$ from Fourier inversion cancels the
  $c_d^{-1}$ in the ordinary Fourier integral of $\psi$. The norms of
  $H^r(\R^d)$ and $H^r(\mathbb{T}_L)$ carry the weights $\langle\xi\rangle^r$
  and $\langle\kappa_k\rangle^r$. After weighting, that is, as a map from
  $\langle\xi\rangle^r\widehat{u}(\xi)$ to
  $\langle\kappa_k\rangle^r(\mathcal{P}_\psi u)_k$, its exact kernel is
  \[
    L^{-d/2}\left(\frac{\langle\kappa_k\rangle}{\langle\xi\rangle}\right)^r
    \widehat{\psi}(\kappa_k-\xi).
  \]
  We turn to $\mathcal{E}_\chi$. For a finite torus polynomial $v$, the
  normalized Fourier kernel of $\mathcal{E}_\chi v$ is
  \[
    L^{-d/2}\left(\frac{\langle\xi\rangle}{\langle\kappa_k\rangle}\right)^r
    \widehat{\chi}(\xi-\kappa_k).
  \]
  There is no additional $c_d$ in this kernel.

  We now bound both kernels. By~\eqref{E:cutoff-peetre}, each is bounded in
  absolute value by a rapidly decreasing function of its frequency difference,
  times a fixed constant. For the first kernel this bound is
  \[
    L^{-d/2}2^{|r|/2}\langle\kappa_k-\xi\rangle^{|r|}
    \left|\widehat{\psi}(\kappa_k-\xi)\right|,
  \]
  and for the second kernel it is the same expression with
  $\widehat{\chi}(\xi-\kappa_k)$ in place of $\widehat{\psi}(\kappa_k-\xi)$.
  These functions of the frequency difference decrease rapidly by the second
  fact of Step~1. One Schur quantity is its $L^1$ integral and the other a
  uniformly translated lattice sum, bounded by~\eqref{E:cutoff-lattice-sum}.
  Indeed, the integral over $\xi$ at fixed $k$ is the $L^1$ integral of the
  bounding function and does not depend on $k$. The sum over $k$ at fixed
  $\xi$ is a lattice sum with spacing $a = 2\pi/L$ and shift $y = \xi$. We
  bound the rapidly decreasing function by a constant times
  $\langle\cdot\rangle^{-N}$ for some $N > d$,
  and~\eqref{E:cutoff-lattice-sum} bounds the sum uniformly in $\xi$. This
  proves~\eqref{E:cutoff-periodic-maps}, including negative $r$.

  \medskip
  \noindent\textit{Step 4.~}We extend the two maps by density. We check that
  the two classes used in
  Step~3 are dense. Smooth compact approximation in weighted Fourier $L^2$
  proves Schwartz density on $\R^d$. Weighted sequence truncation proves
  polynomial density on the torus. Hence, these maps extend uniquely. Their
  distributional meanings follow by testing against compact smooth functions
  and density: each periodic coefficient pairs the distribution with a compact
  smooth test. No point sampling of a rough Fourier function is used. The same
  full-Lebesgue Schur argument, with Lebesgue measure on both sides, proves
  smooth compact multiplication bounded on every whole-space $H^r$.

  \medskip
  \noindent\textit{Step 5.~}We prove the
  identity~\eqref{E:cutoff-whole-space-recovery}. We first
  prove the identity for Schwartz $u$. For Schwartz $u$, the periodic Fourier
  series of $\psi u$ converges smoothly on the cube. Its coefficients are the
  numbers $(\mathcal{P}_\psi u)_k$ of Step~3, and its partial sums are finite
  torus polynomials, on which $\mathcal{E}_\chi$ is multiplication by $\chi$.
  Rapid decay of the Fourier coefficients gives convergence of these partial
  sums in $H^{\sigma+\gamma}(\mathbb{T}_L)$, hence also in
  $H^\sigma(\mathbb{T}_L)$. Boundedness of $\mathcal{E}_\chi$ gives convergence
  of their images to $\mathcal{E}_\chi\mathcal{I}_\gamma\mathcal{P}_\psi u$
  in $H^\sigma(\R^d)$. Their smooth convergence on the support of $\chi$ also
  gives convergence of the products with $\chi$ to $\chi\psi u$ as
  distributions. Since Sobolev convergence implies distributional
  convergence for every real $\sigma$, these two limits agree.
  Since $\chi\psi = \chi$ and $\chi$ is supported strictly inside the cube,
  multiplying this series by $\chi$ recovers $\chi u$ with no alias term.
  Since $\mathcal{I}_\gamma$ is the identity inclusion, this
  is~\eqref{E:cutoff-whole-space-recovery} for Schwartz $u$.

  It remains to pass to all of $H^{\sigma+\gamma}(\R^d)$. Both sides
  of~\eqref{E:cutoff-whole-space-recovery} are bounded from
  $H^{\sigma+\gamma}(\R^d)$ to $H^\sigma(\R^d)$, using the just-proved maps,
  the inclusion and smooth multiplication. In detail, the left side uses
  $\mathcal{P}_\psi$ with $r = \sigma+\gamma$, the inclusion
  $\mathcal{I}_\gamma$, which is bounded because $\gamma > 0$, and
  $\mathcal{E}_\chi$ with $r = \sigma$. The right side uses the multiplication
  by $\chi$ of Step~4 with $r = \sigma$, together with the inclusion of
  $H^{\sigma+\gamma}(\R^d)$ in $H^\sigma(\R^d)$. Since Schwartz functions are
  dense in $H^{\sigma+\gamma}(\R^d)$ by Step~4, density
  proves~\eqref{E:cutoff-whole-space-recovery} on that full domain. This
  proves Lemma~\ref{L:cutoff-factorization}.
\end{proof}

\begin{proposition}[Whole-space cutoff singular values]\label{P:cutoff-singular-values}
  For real $\sigma$ and $\gamma > 0$, the initially Schwartz-defined operator
  $K_{\chi,\sigma,\gamma} = \Lambda^\sigma\chi\Lambda^{-\sigma-\gamma}$ extends
  uniquely to a compact operator on $L^2(\R^d)$, and
  \begin{equation}\label{E:cutoff-singular-values}
    s_j(K_{\chi,\sigma,\gamma})
    \le C_{\chi,\sigma,\gamma,d}j^{-\gamma/d},\qquad j \ge 1.
  \end{equation}
  Consequently, $K_{\chi,\sigma,\gamma} \in S_p$ whenever $p \ge 1$ and $p\gamma
  > d$.
\end{proposition}

\begin{proof}
  We use the cutoff factorization to approximate the operator by finite-rank
  operators. Their ranks and approximation errors will give the singular-value
  bound, which we then sum to obtain Schatten membership.

  \medskip
  \noindent\textit{Step 1.~}We construct finite-rank approximations.
  The zero cutoff is immediate. Fix one cube and cutoff pair $L,\psi$ as in
  Lemma~\ref{L:cutoff-factorization}. Its recovery identity gives
  \[
    K_{\chi,\sigma,\gamma}
    =\Lambda^\sigma\mathcal{E}_\chi\mathcal{I}_\gamma
    \mathcal{P}_\psi\Lambda^{-\sigma-\gamma}.
  \]
  The cutoff maps are bounded by the preceding lemma. Thus, the only
  factor we need to approximate is the torus inclusion.
  In the weighted torus coordinates, $\mathcal{I}_\gamma$ is diagonal with
  entries $\langle\kappa_k\rangle^{-\gamma}$. Let $Q_n$ retain $|k|_\infty \le
  n$. Its rank is $(2n+1)^d$, and
  \[
    \|(1-Q_n)\mathcal{I}_\gamma\|
    \le [1+(2\pi/L)^2(n+1)^2]^{-\gamma/2}\longrightarrow0.
  \]
  Replacing the inclusion by $Q_n\mathcal{I}_\gamma$ gives finite-rank $K_n$
  converging in norm to $K_{\chi,\sigma,\gamma}$. This proves compactness:
  finite-dimensional images have finite nets of each positive radius, and norm
  approximation transfers those nets to the limiting image of the unit ball.

  \medskip
  \noindent\textit{Step 2.~}We compare approximation rank with
  singular-value index.
  Put
  \[
    M=\|\mathcal{E}_\chi:H^\sigma(\mathbb{T}_L)\to H^\sigma(\R^d)\|
    \|\mathcal{P}_\psi:H^{\sigma+\gamma}(\R^d)
    \to H^{\sigma+\gamma}(\mathbb{T}_L)\|.
  \]
  The outer Bessel maps are unitary. For $j \ge 3^d$, choose $n = \lfloor
  j^{1/d}/4\rfloor$. Then
  \[
    (2n+1)^d\le(j^{1/d}/2+1)^d<j,\qquad
    \|K-K_n\|\le M(2L/\pi)^\gamma j^{-\gamma/d}.
  \]
  The strict rank inequality lets us use $K_n$ to bound the $j$th singular
  value. The finitely many smaller indices will be treated by the operator
  norm.
  The second estimate follows from $n+1 > j^{1/d}/4$. For a compact operator,
  every rank less than $j$ approximation $F$ satisfies $s_j(K) \le \|K-F\|$: the
  span of the first $j$ right singular vectors contains a unit vector in $\ker
  F$, on which $\|K\cdot\| \ge s_j(K)$. If $s_j(K) = 0$ the inequality is
  immediate. For $j < 3^d$ use $s_j(K) \le \|K\| \le M$. Thus, one may take
  \[
    C_{\chi,\sigma,\gamma,d}
    =M\max(3^\gamma,(2L/\pi)^\gamma).
  \]
  \medskip
  \noindent\textit{Step 3.~}We sum the singular-value bound.
  Finally, $\sum_j j^{-p\gamma/d}$ converges for $p\gamma > d$. This proves
  Proposition~\ref{P:cutoff-singular-values}.
\end{proof}

No uniform support-radius constant or endpoint necessity is asserted in
Proposition~\ref{P:cutoff-singular-values}. The periodic inclusion was only a
factor of the recovered whole-space operator.

\subsection{Balanced adjoints and localized variance envelopes}

For $E = H^{1/2}(\R^d)\times H^{-1/2}(\R^d)$ use the physical unitary
coordinates
\[
  J_0(u,v)=(\Lambda^{1/2}u,\Lambda^{-1/2}v),\qquad
  H=L^2(\R^d;\mathbb{C}^2).
\]
The frequency map $J$ in Section~\ref{S:wave-variance} is
$(\mathcal{F}\oplus\mathcal{F})J_0$. Let $D_\chi$ multiply both phase entries
and define
\begin{align}
  P_\chi&=J_0D_\chi J_0^{-1},\qquad
  D_\gamma=\operatorname{diag}(\Lambda^{-\gamma},\Lambda^{-\gamma}),
  \label{E:compact-balanced-cutoff}\\
  K_{\rm out}&=P_\chi D_\gamma,\qquad
  K_{\rm in}=P_\chi^*D_\gamma.
  \label{E:compact-two-factors}
\end{align}

\begin{proposition}[Both compact envelopes]\label{P:compact-two-envelopes}
  The operator $P_\chi$ is bounded. For the unitary swap $U(v,w) = (w,v)$,
  \[
    P_\chi^*=UP_\chi U,\qquad K_{\rm in}=UK_{\rm out}U.
  \]
  Both factors are compact and have equal $S_p$ norms for $p \ge 1$, $p\gamma >
  d$. For $p \ge 2$, $p\gamma > d$, the positive compact operators
  \[
    B_{\rm out}=P_\chi D_\gamma^2P_\chi^*=K_{\rm out}K_{\rm out}^*,
    \qquad B_{\rm in}=P_\chi^*D_\gamma^2P_\chi=K_{\rm in}K_{\rm in}^*
  \]
  satisfy
  \begin{equation}\label{E:compact-envelope-traces}
    \operatorname{Tr}B_{\rm out}^{p/2}
    =\operatorname{Tr}B_{\rm in}^{p/2}=\|K_{\rm out}\|_{S_p}^p.
  \end{equation}
\end{proposition}

\begin{proof}
  We first identify the adjoint of the balanced cutoff. The swap identity
  then reduces both factors to the same singular values, and these values
  give the trace identities.

  \medskip
  \noindent\textit{Step 1.~}We identify the adjoint.
  Put $A_\sigma = \Lambda^\sigma\chi\Lambda^{-\sigma}$. The full-Lebesgue Schur
  argument in Lemma~\ref{L:cutoff-factorization} shows boundedness for every
  real $\sigma$. On Schwartz functions, self-adjointness of real multiplication
  and of the Bessel powers gives $A_\sigma^* = A_{-\sigma}$. Density and
  boundedness extend this to the Hilbert adjoints. No evenness of $\chi$ is
  needed. Thus,
  \[
    P_\chi=\operatorname{diag}(A_{1/2},A_{-1/2}),\qquad
    P_\chi^*=\operatorname{diag}(A_{-1/2},A_{1/2})=UP_\chi U.
  \]
  \medskip
  \noindent\textit{Step 2.~}We obtain compactness and Schatten membership.
  The swap commutes with $D_\gamma$, giving the claimed factor identity. The
  blocks of $K_{\rm out}$ are $K_{\chi,1/2,\gamma}$ and $K_{\chi,-1/2,\gamma}$.
  Proposition~\ref{P:cutoff-singular-values} gives compactness and $S_p$
  membership. The singular values of their direct sum form the multiset union,
  so
  \[
    \|K_{\rm out}\|_{S_p}^p
    =\|K_{\chi,1/2,\gamma}\|_{S_p}^p+
    \|K_{\chi,-1/2,\gamma}\|_{S_p}^p.
  \]
  Unitary conjugation preserves them.

  \medskip
  \noindent\textit{Step 3.~}We compute the two traces.
  Finally, the positive eigenvalues of
  $KK^*$ and $K^*K$ agree with multiplicity: on an eigenspace with eigenvalue
  $\lambda > 0$, the maps $K/\sqrt\lambda$ and $K^*/\sqrt\lambda$ are inverse
  isometries. Zero eigenvalues contribute nothing to the trace. Summing these
  scalar powers proves~\eqref{E:compact-envelope-traces}. This proves
  Proposition~\ref{P:compact-two-envelopes}.
\end{proof}

We use the usual positive square root of a bounded positive operator and the
spectral decomposition of compact positive operators. The needed order argument
is recorded explicitly.

\begin{lemma}[Compact positive order and trace]\label{L:compact-positive-order}
  If $A,B$ are bounded self-adjoint operators with $0 \le A \le B$ and $B$
  compact positive, then $A$ is compact and its ordered eigenvalues satisfy
  $\lambda_j(A) \le \lambda_j(B)$, including zero padding. Consequently, for $t
  \ge 1$ with finite $\operatorname{Tr}B^t$,
  \begin{equation}\label{E:compact-positive-trace}
    \operatorname{Tr}A^t\le\operatorname{Tr}B^t.
  \end{equation}
\end{lemma}

\begin{proof}
  Let $Q_n$ be finite spectral projections of $B$ whose tails tend to zero in
  norm. Then
  \[
    \|(1-Q_n)A^{1/2}\|^2
    =\|(1-Q_n)A(1-Q_n)\|
    \le\|(1-Q_n)B(1-Q_n)\|\longrightarrow0.
  \]
  The finite-rank approximants here are $Q_nA^{1/2}$; the displayed norm
  is precisely their error.
  Thus, $A^{1/2}$ is a norm limit of finite-rank operators and is compact, as is
  $A$. For compact positive operators, the max-min formula for $\lambda_j$
  follows in an eigenbasis: the span of the first $j$ eigenvectors gives the
  lower bound for the supremum over $j$-dimensional subspaces of their minimum
  Rayleigh quotient, while every such subspace meets the orthogonal complement
  of the first $j-1$ eigenvectors, giving the upper bound. Compare the Rayleigh
  quotients of $A$ and $B$ to obtain the eigenvalue inequality. Sum the ordered
  scalar powers to prove \eqref{E:compact-positive-trace}. No inequality $A^t
  \le B^t$ is asserted. This proves Lemma~\ref{L:compact-positive-order}.
\end{proof}

\begin{proposition}[Localized finite variance bounds]\label{P:localized-finite-variance}
  Suppose finitely many bounded operators $Z_\ell$ on $H$ satisfy
  \begin{equation}\label{E:localized-finite-hypotheses}
    P_\chi Z_\ell=Z_\ell=Z_\ell P_\chi,\qquad
    V_{\rm out}:=\sum_\ell Z_\ell Z_\ell^*\le CD_\gamma^2,\qquad
    V_{\rm in}:=\sum_\ell Z_\ell^*Z_\ell\le CD_\gamma^2,
  \end{equation}
  with $C \ge 0$. If $p \ge 2$ and $p\gamma > d$, every $Z_\ell$ is compact in
  $S_p$, the two variances are compact in $S_{p/2}$, and
  \begin{align}
    V_{\rm out}&\le CB_{\rm out},\qquad V_{\rm in}\le CB_{\rm in},
    \label{E:localized-two-envelopes}\\
    \operatorname{Tr}V_{\rm out}^{p/2},\quad
    \operatorname{Tr}V_{\rm in}^{p/2}
    &\le C^{p/2}\|K_{\rm out}\|_{S_p}^p.
    \label{E:localized-trace-bounds}
  \end{align}
\end{proposition}

\begin{proof}
  We first use the localization identities to replace the two variance
  bounds by compact envelopes. We then apply the same positive-order argument
  to each individual operator.

  \medskip
  \noindent\textit{Step 1.~}We bound the two variances.
  Left localization gives $V_{\rm out} = P_\chi V_{\rm out}P_\chi^*$. Apply
  positive congruence to the output bound
  in~\eqref{E:localized-finite-hypotheses} to get $V_{\rm out} \le CP_\chi
  D_\gamma^2P_\chi^* = CB_{\rm out}$. Right localization similarly gives $V_{\rm
  in} = P_\chi^*V_{\rm in}P_\chi \le CB_{\rm in}$.
  Positive congruence preserves order: one tests the original inequality
  on the adjoint of the multiplying operator applied to each vector.
  Proposition~\ref{P:compact-two-envelopes} and
  Lemma~\ref{L:compact-positive-order} prove compactness and the two trace
  bounds with $t = p/2$. If $C = 0$, both variances and every $Z_\ell$ vanish,
  so all assertions are immediate.

  \medskip
  \noindent\textit{Step 2.~}We obtain compactness of each summand.
  For each $\ell$, $Z_\ell Z_\ell^* \le V_{\rm out}$ is compact by the same
  lemma. Its finite spectral projections $Q_n$ give $\|(1-Q_n)Z_\ell\|^2 =
  \|(1-Q_n)Z_\ell Z_\ell^*(1-Q_n)\| \to 0$. Hence, $Z_\ell$ is compact. Its
  singular values squared are the nonzero eigenvalues of $Z_\ell Z_\ell^*$.
  Comparing with $V_{\rm out}$ and summing proves $Z_\ell \in S_p$ with
  $\|Z_\ell\|_{S_p}^p \le \operatorname{Tr}V_{\rm out}^{p/2}$. This proves
  Proposition~\ref{P:localized-finite-variance}.
\end{proof}

The two localization identities are hypotheses here; a noise coefficient family
has not yet been shown to satisfy them. In particular, $P_\chi$ is not assumed
self-adjoint or a projection.

\begin{corollary}[Finite Gaussian operators]\label{C:finite-gaussian-operators}
  Under~\eqref{E:localized-finite-hypotheses}, let $g_\ell$ be independent real
  standard Gaussians. For even $p \ge 2$ with $p\gamma > d$, the finite sum $X =
  \sum_\ell g_\ell Z_\ell$ satisfies
  \begin{equation}\label{E:finite-gaussian-operator-bound}
    (\E\|X\|_{S_p}^p)^{1/p}
    \le\sqrt{p-1}\sqrt C\|K_{\rm out}\|_{S_p}.
  \end{equation}
  Its extended norm is measurable, with finite expectation as displayed.
\end{corollary}

\begin{proof}
  We compress the operators to finite matrices, where the Gaussian moment
  bound is already available. We then pass to the original operators using
  convergence of singular values and two applications of Fatou's lemma.

  \medskip
  \noindent\textit{Step 1.~}We bound the variances of the compressions.
  Choose finite-rank orthogonal projections $\Pi_n$ from a countable orthonormal
  basis of $H$, with $\Pi_n \to I$ strongly. Put $Z_{\ell,n} = \Pi_nZ_\ell\Pi_n$
  and $X_n = \sum_\ell g_\ell Z_{\ell,n}$, viewed as matrices on
  $\operatorname{ran}\Pi_n$. Their output variance satisfies
  \[
    \sum_\ell\Pi_nZ_\ell\Pi_nZ_\ell^*\Pi_n
    \le\Pi_nV_{\rm out}\Pi_n,
  \]
  since the difference is $\sum_\ell\Pi_nZ_\ell(1-\Pi_n)Z_\ell^*\Pi_n \ge 0$.
  The input inequality is identical with adjoints interchanged. Compact max-min
  also gives $\lambda_j(\Pi_nV\Pi_n) \le \lambda_j(V)$ for positive compact $V$,
  since Rayleigh subspaces inside $\operatorname{ran}\Pi_n$ form a subset of all
  subspaces. Thus, both variance trace powers of these matrices are bounded by
  $C^{p/2}\|K_{\rm out}\|_{S_p}^p$.

  \medskip
  \noindent\textit{Step 2.~}We apply the finite matrix estimate.
  Proposition~\ref{P:finite-gaussian-matrix} applies to these finite complex
  matrices, the same real independent Gaussians and the even $p$. Taking the
  average of the two already bounded trace powers gives
  \begin{equation}\label{E:finite-gaussian-compression-bound}
    \E\|X_n\|_{S_p}^p
    \le(p-1)^{p/2}C^{p/2}\|K_{\rm out}\|_{S_p}^p.
  \end{equation}
  No factor two is added by this average.

  \medskip
  \noindent\textit{Step 3.~}We pass to the limit and check measurability.
  This passage is pointwise in the Gaussian coordinates until the final
  application of Fatou in probability.
  For each Gaussian coordinate vector, $X$ is a finite sum of compact operators.
  Strong convergence of $\Pi_n$ implies $\Pi_nX\Pi_n \to X$ in operator norm:
  approximate $X$ by a finite-rank $F$, use uniform strong convergence on the
  finite-dimensional ranges of $F$ and $F^*$, and control the remainder using
  $\|\Pi_n\| \le 1$. Singular values are norm-Lipschitz. Indeed, their formula
  \[
    s_j(A)=\inf_{\dim L<j}
    \sup_{v\in L^\perp,\ \|v\|=1}\|Av\|
  \]
  follows from compact singular vectors and gives $s_j(A) \le s_j(B)+\|A-B\|$,
  and the reverse bound by interchanging $A,B$. Hence, $s_j(X_n) \to s_j(X)$ for
  each $j$. Fatou for the nonnegative singular-value sum gives $\|X\|_{S_p}^p
  \le \liminf_n\|X_n\|_{S_p}^p$, allowing extended infinity. These variables are
  measurable: the finite Gaussian coordinate map into operators is
  norm-continuous, each singular value is continuous, and the norm power is
  their nonnegative measurable sum. Fatou in probability
  and~\eqref{E:finite-gaussian-compression-bound} prove the bound and
  almost-sure finiteness. No $S_p$-Bochner convergence is assumed. This proves
  Corollary~\ref{C:finite-gaussian-operators}.
\end{proof}

\subsection{A nonzero witness and the need for a compact cutoff}

Proposition~\ref{P:localized-finite-variance} and
Corollary~\ref{C:finite-gaussian-operators} assume
\eqref{E:localized-finite-hypotheses}. We check that these hypotheses admit a
nonzero operator, and then test what is lost without the compact cutoff.

\begin{proposition}[Whole-space cutoff tests]\label{P:cutoff-boundary-tests}
  If $\chi$ equals one on a nonempty open set, there is a nonzero rank-one
  operator satisfying~\eqref{E:localized-finite-hypotheses} with a finite common
  $C$. On the other hand, $\Lambda^{-\gamma}$ on $L^2(\R^d)$ is not compact for
  any $\gamma > 0$.
\end{proposition}

\begin{proof}
  We treat the two assertions separately. A rank-one operator built from test
  functions supported on the plateau provides the witness. Translates of a
  fixed function show why smoothing alone does not give compactness on the
  whole space.

  \medskip
  \noindent\textit{Step 1.~}We construct the localized witness.
  Choose nonzero real $v,f \in C_c^\infty$ supported inside the plateau, and put
  $\widetilde{v} = (\Lambda^{1/2}v,0)$ and $\widetilde{w} =
  (\Lambda^{-1/2}f,0)$. Both are Schwartz. The block formulas give
  $P_\chi\widetilde{v} = \widetilde{v}$ and $P_\chi^*\widetilde{w} =
  \widetilde{w}$. Therefore, $Zx = \langle x,\widetilde{w}\rangle\widetilde{v}$
  satisfies $P_\chi Z = Z = ZP_\chi$. Its two variances are
  $\|\widetilde{w}\|^2\widetilde{v}\otimes\widetilde{v}^*$ and
  $\|\widetilde{v}\|^2\widetilde{w}\otimes\widetilde{w}^*$. For any Schwartz $a$
  and any $x \in H$,
  \[
    |\langle x,a\rangle|
    =|\langle D_\gamma x,D_\gamma^{-1}a\rangle|
    \le\|D_\gamma x\|\|D_\gamma^{-1}a\|.
  \]
  The last inequality is Cauchy--Schwarz. The preceding equality uses
  self-adjointness of the multiplier and its inverse on Schwartz functions.
  The inverse multiplier is applied only to Schwartz $a$, so it belongs to $H$.
  Apply the inequality with $a = \widetilde{v},\widetilde{w}$ to obtain the two
  variance bounds with
  \[
    C=\max\left(
    \|\widetilde{w}\|^2\|D_\gamma^{-1}\widetilde{v}\|^2,
    \|\widetilde{v}\|^2\|D_\gamma^{-1}\widetilde{w}\|^2\right)<\infty.
  \]
  This witness uses a plateau; no nonzero family is inferred for every compact
  cutoff or identified with a source family.

  \medskip
  \noindent\textit{Step 2.~}We test compactness by translation.
  For noncompactness, take $f$ with nonzero Gaussian Fourier transform. Then $g
  = \Lambda^{-\gamma}f$ is a nonzero Schwartz function. The images of translates
  of $f$ are the same translates of $g$. With $T_hg(x) = g(x-h)$, Plancherel
  gives
  \[
    \langle T_hg,g\rangle
    =\int e^{-i\xi\cdot h}|\widehat{g}(\xi)|^2\ud\xi
    =c_d^{-1}\widehat{|\widehat{g}|^2}(h).
  \]
  The Fourier multiplier commutes with translations, so the input
  sequence and its image have the same translation parameters.
  This tends to zero because $|\widehat{g}|^2$ is Schwartz. Choose a large
  spacing so that any two distinct translates in the resulting sequence have
  image distance at least $\|g\|$. The input norms are identical, but the images
  have no Cauchy subsequence. This disproves compactness. This proves
  Proposition~\ref{P:cutoff-boundary-tests}.
\end{proof}

The translation obstruction is restated in No-go~\ref{P:nogo-24}:
frequency decay alone does not replace compact spatial localization.


Thus, the choice $\chi = 1$ cannot be substituted into the compact-cutoff
argument. No necessity at the Schatten endpoint $p\gamma = d$ is asserted. The
results concern finite matrices and conditional finite operator sums.
Coefficient localization, compact variance tails, an infinite Gaussian series
and time regularity are proved in Section~\ref{S:actual-operator-driver}; source
and Young identification are proved in
Sections~\ref{S:source-driver-convergence}--\ref{S:constant-phase-young}. They
are not consequences of the finite comparison alone.
